\begingroup
\small
\renewcommand{\arraystretch}{1.08}
\setlength{\LTleft}{0pt}
\setlength{\LTright}{0pt}

\begin{longtable}{|p{0.28\textwidth}|p{0.66\textwidth}|}
\hline
\textbf{Термин} & \textbf{Определение} \\
\hline
\endfirsthead

\hline
\textbf{Термин} & \textbf{Определение} \\
\hline
\endhead

Apache Arrow & Открытая платформа для представления и обмена колонoчными данными в оперативной памяти~\cite{apache-arrow}. \\
\hline
DuckDB & Встраиваемая аналитическая СУБД, ориентированная на обработку данных в колонoчном представлении~\cite{duckdb-extensions}. \\
\hline
Lance & Колонoчный формат и формат набора данных, ориентированный на аналитическую обработку и быструю выборку~\cite{lance-format}. \\
\hline
ORC & Колонoчный формат хранения данных, применяемый в аналитических системах и ориентированный на эффективное сжатие и сканирование~\cite{apache-orc}. \\
\hline
Parquet & Колонoчный формат хранения данных, предназначенный для аналитической обработки и межсистемного обмена~\cite{apache-parquet}. \\
\hline
Scale factor & Параметр спецификации TPC-H, определяющий размер формируемого набора данных в гигабайтах~\cite{tpch-spec}. \\
\hline
TPC-H & Стандартный аналитический бенчмарк, включающий фиксированную схему таблиц, правила генерации данных и набор запросов для оценки производительности СУБД~\cite{tpch-spec}. \\
\hline
Vortex & Адаптивный колонoчный формат данных для систем, использующих Apache Arrow~\cite{vortex-format}. \\
\hline
Lakehouse-система & Класс аналитических систем, объединяющих хранение данных в файловом озере с табличными слоями метаданных; примером такого слоя является Apache Iceberg~\cite{apache-iceberg}. \\
\hline
OLAP & Класс систем и методов аналитической обработки данных, ориентированных на выполнение сложных запросов над большими массивами информации. \\
\hline
Amazon S3 & Объектное хранилище и программный интерфейс доступа к данным, де-факто используемые как стандарт для совместимых облачных хранилищ. \\
\hline
Батч & Группа строк, объединенная для совместной обработки, передачи или сериализации. \\
\hline
Партиция & Логическая часть таблицы, генерируемая и обрабатываемая как самостоятельный фрагмент общего набора данных. \\
\hline
Сериализация & Преобразование внутреннего представления данных в формат, пригодный для записи на носитель или передачи между компонентами. \\
\hline
СУБД & Система управления базами данных, обеспечивающая хранение, обработку и доступ к данным. \\
\hline
\end{longtable}

\normalsize
\endgroup
